% 出处：主技能 skills/math-modeling-helper/SKILL.md「LaTeX模板（精简版）」小节（原样提取，占位示例保留）
% 用法：复制到工作区 paper/paper.tex 填写内容 → xelatex 编译两遍 → pandoc paper.tex -o paper.docx → python code/word_postprocess.py paper/paper.docx 版式微调
\documentclass[12pt,a4paper]{article}

% ========== 编码与中文支持 ==========
\usepackage{ctex}                    % 中文支持（禁止inputenc）

% ========== 宏包 ==========
\usepackage{amsmath,amssymb,amsfonts}
\usepackage{bm}                      % 矩阵/向量加粗
\usepackage{geometry}
\usepackage{graphicx}
\usepackage{float}
\usepackage{booktabs}                % 三线表
\usepackage{array}                   % m{}列类型（垂直居中）
\usepackage{tabularx}                % X列类型（自动宽度换行）
\usepackage{multirow}
\usepackage{caption}
\usepackage{subcaption}
\usepackage{hyperref}
\usepackage{fancyhdr}
\usepackage{titlesec}
\usepackage{listings}
\usepackage{xcolor}
\usepackage{seqsplit}                % 长单词断行

% ========== 西文字体 Times New Roman ==========
\setmainfont{Times New Roman}        % 西文和数字使用 Times New Roman

% ========== 页面设置（四边2.5cm） ==========
\geometry{top=2.5cm, bottom=2.5cm, left=2.5cm, right=2.5cm}

% ========== 标题格式（2026规范） ==========
\renewcommand{\thesection}{\chinese{section}}  % 一级标题使用中文数字
\titleformat{\section}{\centering\heiti\fontsize{14}{18}\selectfont}{\thesection、}{1em}{}  % 标题显示加顿号，交叉引用不受影响
\titleformat{\subsection}{\heiti\fontsize{12}{16}\selectfont}{\arabic{section}.\arabic{subsection}}{1em}{}
\titleformat{\subsubsection}{\heiti\fontsize{12}{16}\selectfont}{\arabic{section}.\arabic{subsection}.\arabic{subsubsection}}{1em}{}

% ========== 正文格式 ==========
\setlength{\parindent}{2em}
\setlength{\parskip}{0pt}
\renewcommand{\baselinestretch}{1}  % 1倍行距

% ========== 图表标题 ==========
\captionsetup{font=small, labelsep=space, skip=6pt}
\captionsetup[figure]{position=below}
\captionsetup[table]{position=above}

% ========== 页眉页脚：禁止页眉；页码从第一页（摘要页）起页尾居中 ==========
\pagestyle{fancy}
\fancyhf{}                        % 清空页眉页脚（页眉必须留空）
\fancyfoot[C]{\thepage}           % 页尾居中页码，从摘要页第1页开始
\renewcommand{\headrulewidth}{0pt} % 禁止页眉横线（保证无页眉）

% ========== 代码样式 ==========
\lstset{
    basicstyle=\small\ttfamily,
    keywordstyle=\color{blue},
    commentstyle=\color{green!60!black},
    stringstyle=\color{red!70!black},
    numbers=left,
    numberstyle=\tiny\color{gray},
    frame=single,
    breaklines=true
}

\begin{document}

% ========== 第一页：题目 + 摘要 + 关键词（必须同页，页码从本页开始） ==========
\title{\heiti\fontsize{16}{20}\selectfont 论文题目}
\author{}
\date{}
\maketitle
\thispagestyle{fancy}  % 覆盖 \maketitle 默认的 plain 样式，确保首页页脚与正文一致

\vspace{0.3cm}

% 摘要（与标题、关键词同一页，手动格式与正文一致）
% 按五段式组织：① 钩子切入（背景中的具体矛盾/关键数字/决策困境，最多两句）② 方法链与一句话创新点
% ③ 逐问结果（模型+算法+关键数值+对比）④ 稳健性/误差一句 ⑤ 结论的现实映射；
% 禁止"随着…的发展""本文对…进行了深入研究""得到了较为理想的结果"等模板腔；详见 abstract-template.md
\noindent{\heiti\fontsize{14}{18}\selectfont 摘要：}\par
\setlength{\parindent}{2em}
\normalsize
[① 钩子：以背景中的具体矛盾或关键数字直入……② 方法链与创新点：本文以…为主线，建立…模型，并针对…设计…算法……③ 逐问结果：问题一…[关键数值]，较…[对比]；问题二以该基线为参照…；问题三进一步……④ 稳健性：敏感性分析显示…结论稳健……⑤ 映射：所得方案/结论可为[具体场景]提供决策依据]

\vspace{0.2cm}

\noindent{\heiti 关键词：}关键词1；关键词2；关键词3

% ========== 关键词后换页，正文从第2页开始 ==========
\newpage

% ========== 正文内容 ==========
% 一、问题背景与重述
% 二、问题分析
% 三、模型假设
\section{模型假设}

针对本题实际情况，在建模前提出如下假设，如表\ref{tab:assumptions}所示。

\begin{table}[H]
\centering
\caption{模型假设}
\label{tab:assumptions}
\begin{tabularx}{\textwidth}{c X X}
\toprule
编号 & \multicolumn{1}{c}{假设内容} & \multicolumn{1}{c}{合理性说明} \\
\midrule
1 & 假设数据在采集过程中无系统性偏差 & 依据题目说明，数据经标准化预处理 \\
2 & 假设各周期决策相互独立，无跨周期资源传递 & 题目未声明资源可累积，故按独立处理 \\
3 & 假设模型参数在求解周期内保持恒定 & 短时间尺度内参数漂移可忽略\cite{ref1} \\
4 & 假设随机扰动服从均值为零的正态分布 & 依据中心极限定理，多源扰动叠加趋近正态 \\
\bottomrule
\end{tabularx}
\end{table}

% 四、符号说明
% 五、模型建立与求解（若所有问题共用一套数据，则先设"五、数据处理"，本章顺延为第六章，后续章节类推）
% 六、模型的评价、改进与推广

% ========== 参考文献（另起一页，页码连续） ==========
\newpage
\begin{thebibliography}{99}
\bibitem{ref1} ...
\end{thebibliography}

% ========== 附录/支撑材料（另起一页，页码连续） ==========
\newpage
\section*{附录}
% 附录A 支撑材料清单
% 附录B 核心代码

\end{document}
